% !TEX program = lualatex
% !TEX encoding = UTF-8
% !BIB program = bibtex
% !TEX spellcheck = en_US
\documentclass[12pt]{article}
\usepackage{url}
\usepackage{amsmath}
\usepackage{graphicx}
\usepackage{subfigure}
\graphicspath{{images/}}
\usepackage{fancyhdr}
\usepackage{enumerate}
\usepackage[letterpaper, margin=1in]{geometry}
\usepackage{siunitx}
\usepackage{hyperref}
\usepackage{bm}
\usepackage{braket}
\usepackage[version=4]{mhchem}
\usepackage[symbol]{footmisc}  % footnote symbol
\usepackage{relsize}

\sisetup{inter-unit-product=\ensuremath{{}\cdot{}}}
\DeclareTextCommand{\nobreakspace}{T1}{\leavevmode\nobreak\ }

\title{Solving \ce{He} and \ce{H2} within Hartree--Fock}
\author{}
\date{}

\makeatletter
\let\thetitle\@title
\let\theauthor\@author
\let\thedate\@date
\makeatother

\renewcommand{\sectionmark}[1]{\markright{#1}}

\pagestyle{fancy}
\fancyhf{}
\lhead{\fancyplain{}{\rightmark }}
\cfoot{\thepage}
\setlength\headheight{14.5pt}

\newcommand*{\rmn}[1]{\romannumeral#1}
\newcommand*{\RMN}[1]{\uppercase\expandafter{\romannumeral#1}}
\newcommand*{\tran}{^{\mkern-1.5mu\mathsf{T}}}
\renewcommand{\thefootnote}{\fnsymbol{footnote}}

\begin{document}

\section{Problem \RMN{1}}

\subsection*{(a)}

The Hamiltonian is $H = - \frac{1}{2} \nabla^2 -
    \frac{1}{r}$\thinspace,
construct wave functions with $4$ bases,
\begin{equation}
    \psi(\bm{r}) = \sum_{i} c_i \phi_i = c_1 e^{- 13.0 r^2} +
    c_2 e^{- 1.96 r^2} + c_3 e^{- 0.44 r^2} + c_4 e^{-0.12 r^2}.
\end{equation}
where $c_i, i=1,2,3,4$ are real numbers.
Define functional
\begin{equation}
    E[\psi] = \frac{\braket{\psi| H | \psi}}{\braket{\psi|\psi}},
\end{equation}
according to variational principle, the ground state energy $E$ satisfies
\begin{equation}
    \det (H_{ij} - S_{ij} E) = 0,
\end{equation}
where $H_{ij} = \braket{\phi_i | H | \phi_j}$, $S_{ij} =
    \braket{\phi_i | \phi_j}$.
Under spherical coordinates,
\begin{equation}
    \begin{aligned}
        \nabla^2 & = \frac{1}{r} \frac{\partial^2}{\partial r^2} r - \frac{\bm{l}^2}{\hbar^2 r^2} \\
                 & = \frac{1}{r} \frac{\partial^2}{\partial
            r^2}r+{\frac{1}{r^{2}\sin \theta }}{\frac{\partial
        }{\partial \theta }}\Big(\sin \theta {\frac{\partial
            }{\partial \theta }}\Big)+{\frac{1}{r^{2}\sin^{2}\theta }}{\frac{\partial^{2}}{\partial \varphi^{2}}} \thinspace,
    \end{aligned}
\end{equation}
here every $\phi(r)$ does not contain $\theta$ and $\varphi$,
so the Laplacian reduces to
\begin{equation}
    \nabla^2 = \frac{1}{r} \frac{\partial^2}{\partial r^2} r.
\end{equation}
So
\begin{equation}
    \begin{aligned}
        \braket{\phi_i | H | \phi_j} & =
        \int \phi_i^\ast H \phi_j \, d\bm{r} \\
                                     & =
        4\pi\int_{0}^{\infty} \frac{-1 + 3 \alpha_j r - 2 \alpha_j^2 r^3}{r}
        e^{-(\alpha_i + \alpha_j)r^2}
        r^2\, dr.
    \end{aligned}
\end{equation}
Therefore
\begin{equation}
    H =
    \begin{pmatrix}
        0.577368  & 0.0717194 & -0.323206 & -0.437105 \\
        0.0717194 & 0.506469  & -1.00354  & -2.39107  \\
        -0.323206 & -1.00354  & -2.68809  & -7.46151  \\
        -0.437105 & -2.39107  & -7.46151  & -17.6552
    \end{pmatrix},
\end{equation}
and
\begin{equation}
    S =
    \begin{pmatrix}
        0.0420015 & 0.0962338 & 0.113012 & 0.117172 \\
        0.0962338 & 0.717457  & 1.49764  & 1.85622  \\
        0.113012  & 1.49764   & 6.74529  & 13.2875  \\
        0.117172  & 1.85622   & 13.2875  & 47.3596
    \end{pmatrix}.
\end{equation}
Then solve for different set of eigenvalues, we derive Tab. \ref{tab:1-a}.
\begin{table}[h]
    \centering
    \caption{Ground state energy as a function of $n$.}
    \begin{tabular}{|c|c|}
        \hline
        Basis set                                     & Energy (in Hartree) \\
        \hline
        $\{ \alpha_1\}$                               & $13.74637261$       \\
        \hline
        $\{ \alpha_1, \alpha_2\}$                     & $0.705886690023$    \\
        \hline
        $\{ \alpha_1, \alpha_2, \alpha_3\}$           & $-0.434106467896$   \\
        \hline
        $\{ \alpha_1, \alpha_2, \alpha_3, \alpha_4\}$ & $-0.499276143413$   \\
        \hline
    \end{tabular}
    \label{tab:1-a}
\end{table}

Since exact value is $\SI{-13.6}{\electronvolt} \approx \SI{-0.5}{\hartree}$, so the result of
$4$ bases is very close to it, but not with less bases. See Fig. \ref{fig:pro1a} for reference.

% \begin{figure}[h]
%     \centering
%     \includegraphics[width=0.5\linewidth]{pro_1_a}
%     \caption{Ground state energy as a function of $n$.}
%     \label{fig:pro1a}
% \end{figure}

\subsection*{(b)}

For the energy of the first excited state (ie. second lowest Eigenvalue)
as a function of $n$, see Tab. \ref{tab:1-b} and Fig. \ref{fig:pro1b} for reference. Again, the exact value
is $\SI{-3.4}{\electronvolt} \approx \SI{-0.125}{\hartree}$. Now this time even with $4$ bases,
the energy is still far from exact value.
\begin{table}[h]
    \centering
    \caption{Energy of the first excited state as a function of $n$.}
    \begin{tabular}{|c|c|}
        \hline
        Basis set                                     & Energy (in Hartree) \\
        \hline
        $\{ \alpha_1, \alpha_2\}$                     & $19.497200032$      \\
        \hline
        $\{ \alpha_1, \alpha_2, \alpha_3\}$           & $2.15401296835$     \\
        \hline
        $\{ \alpha_1, \alpha_2, \alpha_3, \alpha_4\}$ & $0.106674417746$    \\
        \hline
    \end{tabular}
    \label{tab:1-b}
\end{table}

% \begin{figure}[h]
%     \centering
%     \includegraphics[width=0.5\linewidth]{pro_1_b}
%     \caption{Energy of the first excited state as a function of $n$.}
%     \label{fig:pro1b}
% \end{figure}

\subsection*{(c)}

By solving eigenvalue problem we derive coefficients, that is,
$\begin{pmatrix}
        c_1 & c_2 & c_3 & c_4
    \end{pmatrix}\tran$. The exact ground state is\footnote{Griffith page $138$.}
\begin{equation}
    \psi_{\text{exact}}(\bm{r}) = \frac{1}{\sqrt{\pi}} e^{-r}.  % This function is normalized.
\end{equation}
Then plot as Fig. \ref{fig:pro1c}.
\begin{figure}[h]
    \centering
    \includegraphics[width=0.5\linewidth]{pro_1_c}
    \caption{Ground state eigenfunction as a function of $n$.}
    \label{fig:pro1c}
\end{figure}
From here we can see the wave function with $4$ bases is very close to exact value except $r \approx 0$, but
wave functions with less bases are far from exact value.


\section{Problem \RMN{2}}

\subsection*{(a)}

Set Hamiltonian to be $H = - \frac{1}{2}\nabla_1^2 - \frac{2}{r_1} - \frac{1}{2}\nabla_2^2 - \frac{2}{r_2}$, then
the Schr\"{o}dinger equation is turned into $2$ single-electron equations, \textit{id est},
\begin{equation}
    \Big(- \frac{1}{2}\nabla^2 - \frac{2}{r} \Big) \psi(\bm{r}) = E \psi(\bm{r}).
\end{equation}
Construct single-electron wave function with $4$ bases,
\begin{equation}
    \psi(\bm{r}) = \sum_{i} c_i \phi_i = c_1 e^{- 0.298073 r^2} +
    c_2 e^{- 1.242567 r^2} + c_3 e^{- 5.782948 r^2} + c_4 e^{-38.474970 r^2}.
\end{equation}
where $c_i, i=1,2,3,4$, are real numbers.
Now the ground state energy is \SI{-1.99426616}{\electronvolt}.
The density is
\begin{equation}\label{eq:density}
    \lvert \psi(\bm{r}) \rvert^2= \psi^\ast(\bm{r}) \psi(\bm{r}) =
    \sum_{i, j} c_i^\ast c_j \phi_i^\ast(\bm{r}) \phi_j(\bm{r}),
\end{equation}
where $i, j=1,2,3,4$.
The radial probability distribution is
\begin{equation}
    \lvert \chi(\bm{r}) \rvert^2= r^2\int \lvert \psi(\bm{r}) \rvert^2
    \, d\Omega = 4\pi r^2 \lvert \psi(\bm{r}) \rvert^2.
\end{equation}
And, the exact value of single-electron wave function is
\begin{equation}
    \psi_{\text{exact}}(\bm{r}) = \sqrt{\frac{8}{\pi}} e^{-2 r},  % This function is normalized.
\end{equation}
since the total wave function is
$\Psi_{\text{exact}}(\bm{r}_1, \bm{r}_2) = \frac{8}{\pi}
    e^{-2(r_1 + r_2)}$, see Griffith page $187$ \footnote{Here
    Griffith only considered Hartree method but not
    Hartree--Fock method, so we don't need to construct an
    antisymmetric all-electron wave function.}.
Then we can plot the density and radial probability distribution as a function of $r$
as Fig.  \ref{fig:2-a:mini:subfig:a} and Fig. \ref{fig:2-a:mini:subfig:b}, respectively.

% \begin{figure}[h]
%     \subfigure[Density as a function of $r$.]{
%         \label{fig:2-a:mini:subfig:a}   %% label for first subfigure
%         \begin{minipage}[b]{0.5\textwidth}
%             \centering
%             \includegraphics[width=\linewidth]{pro_2_a_1}
%         \end{minipage}}
%     \hfill
%     \subfigure[Radial probability distribution as a function of $r$.]{
%         \label{fig:2-a:mini:subfig:b}   %% label for second subfigure
%         \begin{minipage}[b]{0.5\textwidth}
%             \centering
%             \includegraphics[width=\linewidth]{pro_2_a_2}
%         \end{minipage}}
%     \caption{Problem $2(a)$.}
%     \label{fig:2-a}   %% label for entire figure
% \end{figure}

\subsection*{(b)}

Now consider the single-electron
Schr\"{o}dinger equation of $\psi_1(\bm{r}_1)$ will be
\begin{equation}
    \Big(- \frac{1}{2}\nabla_1^2 - \frac{2}{r_1}
    + \int \frac{\psi_2^\ast(\bm{r}_2)\psi_2(\bm{r}_2)}{\lvert \bm{r}_1 -
        \bm{r}_2 \rvert} d\bm{r}_2 \Big) \psi_1(\bm{r}_1) = E \psi_1 (\bm{r}_1),
\end{equation}
with $\psi_2(\bm{r}_2) = 0.08704173 e^{-0.298073 r^2}
    + 0.52509737 e^{-1.242567 r^2}
    + 0.51920929 e^{-5.782948 r^2}
    + 0.31233737 e^{-38.474970 r^2}$ as trial function (last question result),
and
\begin{equation}
    \begin{aligned}
        \int \frac{\psi_2^\ast(\bm{r}_2)\psi_2(\bm{r}_2)}{\lvert \bm{r}_1 -
            \bm{r}_2 \rvert} d\bm{r}_2
         & = \sum_{i, j}
        c_i^\ast c_j
        \int
        \frac{\phi_i^\ast(r_2) \phi_j(r_2)}{\lvert \bm{r}_1 - \bm{r}_2 \rvert}
        \, d\bm{r}_2     \\
         & = \sum_{i, j}
        2\pi c_i^\ast c_j
        \int_{0}^{\infty} \int_{0}^{\pi}
        \tfrac{\phi_i^\ast(r_2) \phi_j(r_2)}{\sqrt{r_1^2 + r_2^2 - 2 r_1 r_2
                \cos\theta}}
        r_2^2 \sin \theta \, d\theta \, dr_2.
    \end{aligned}
\end{equation}
Left multiplied by $\bra{\phi_i(\bm{r}_1)}$ and integrate, we have
\begin{equation}
    \begin{aligned}
        \Braket{ \phi_i(\bm{r}_1) | - \frac{1}{2}\nabla_1^2 -
            \frac{2}{r_1} | \psi_1(\bm{r}_1) } +
        \Braket{\phi_i(\bm{r}_1) |
            \int \frac{\psi_2^\ast(\bm{r}_2)\psi_2(\bm{r}_2)}{\lvert
                \bm{r}_1 - \bm{r}_2 \rvert} d\bm{r}_2 | \psi_1(\bm{r}_1)}
        = E \braket{\phi_i(\bm{r}_1) |
            \psi_1(\bm{r}_1)}.
    \end{aligned}
\end{equation}
The eigenvalue problem will be
\begin{equation}
    \det (H_{ij} - S_{ij} E + K_{ij}) = 0,
\end{equation}
here $H_{ij} = \braket{\phi_i | - \frac{1}{2}\nabla_1^2 - \frac{2}{r_1} | \phi_j}$, $S_{ij} =
    \braket{\phi_i | \phi_j}$, and
\begin{equation}
    K_{ij} = \sum_{m, n} 8\pi^2 c_m^\ast c_n
    \int_{0}^{\infty} r_1^2  \, dr_1
    \int_{0}^{\infty} r_2^2 \, dr_2
    \int_{0}^{\pi}
    \tfrac{\phi_i^\ast(r_1) \phi_m^\ast(r_2) \phi_n(r_2)\phi_j(r_1)}{%
        \sqrt{r_1^2 + r_2^2 - 2 r_1 r_2 \cos\theta}} \sin\theta \, d\theta.
\end{equation}
Derive new set of $\{ c_1, c_2, c_3, c_4\}$, let $\psi_1(\bm{r}_1) =
    \sum_{i=1}^{4} c_i \phi_i(\bm{r}_1)$, and substitute
into
\begin{equation}
    \Big(- \frac{1}{2}\nabla_2^2 - \frac{2}{r_2}
    + \int \frac{\psi_1^\ast(\bm{r}_1)\psi_1(\bm{r}_1)}{\lvert \bm{r}_1 -
        \bm{r}_2 \rvert} d\bm{r}_1 \Big) \psi_2(\bm{r}_2) = E \psi_2 (\bm{r}_2),
\end{equation}
repeat this step until $E$ converges.
The total energy is
\begin{equation}
    E_{\text{total}} = 2E -
    \int \frac{\psi_1^\ast(\bm{r}_1)\psi_2^\ast(\bm{r}_2) \psi_1(\bm{r}_1)
        \psi_2(\bm{r}_2)}{\lvert
        \bm{r}_1 - \bm{r}_2 \rvert}\, d\bm{r}_2 \, d\bm{r}_1,
\end{equation}
$2E$ is because $\psi_1 = \psi_2$.
See Fig. \ref{fig:2-b:mini:subfig:a} for reference.

% \begin{figure}[h]
%     \subfigure[The total energy $\varepsilon$ as a function of iteration $n$.]{
%         \label{fig:2-b:mini:subfig:a}   %% label for first subfigure
%         \begin{minipage}[b]{0.5\textwidth}
%             \centering
%             \includegraphics[width=\linewidth]{pro_2_b_1}
%         \end{minipage}}
%     \hfill
%     \subfigure[The interaction density compared with the non-interacting density.]{
%         \label{fig:2-b:mini:subfig:b}   %% label for second subfigure
%         \begin{minipage}[b]{0.5\textwidth}
%             \centering
%             \includegraphics[width=\linewidth]{pro_2_b_2}
%         \end{minipage}}
%     \caption{Problem $2(b)$.}
%     \label{fig:2-b}   %% label for entire figure
% \end{figure}

For density, same as $\eqref{eq:density}$, plot as Fig. \ref{fig:2-b:mini:subfig:b}.

\section{Problem \RMN{3}}

\subsection*{(a)}

By neglecting the electron-electron interactions, we have
\begin{equation}
    H = - \frac{1}{2}\nabla_1^2
    - \frac{1}{2}\nabla_2^2
    - \frac{1}{\lvert \bm{r}_1 - \bm{R}_1 \rvert} - \frac{1}{\lvert \bm{r}_1 - \bm{R}_2 \rvert}
    - \frac{1}{\lvert \bm{r}_2 - \bm{R}_1 \rvert} - \frac{1}{\lvert \bm{r}_2 - \bm{R}_2 \rvert}
    + \frac{1}{\lvert \bm{R}_1 - \bm{R}_2 \rvert},
\end{equation}
The single-electron wave function is
\begin{equation}
    \psi(\bm{r}) = \sum_{i} \sum_n c_{i, m}\phi_{i}^{R_m}(r) =
    \sum_{i=1}^{4}\sum_{m=1}^{2} c_{i, m} e^{-\alpha_i (r - R_m)^2},
\end{equation}
so corresponding Schr\"{o}dinger equation of $\psi_1(\bm{r}_1)$ is
\begin{equation}
    \Big(- \frac{1}{2}\nabla_1^2
    - \frac{1}{\lvert \bm{r}_1 - \bm{R}_1 \rvert}
    - \frac{1}{\lvert \bm{r}_1 - \bm{R}_2 \rvert}
    + \frac{1}{\lvert \bm{R}_1 -
        \bm{R}_2 \rvert} \Big) \psi_1(\bm{r}_1)
    = E \psi_1 (\bm{r}_1).
\end{equation}
Left multiplied by $\Big( \phi_i^{R_m} \Big)^\ast(\bm{r}_1)$ and integrate, we have
\begin{equation}
    \Braket{\phi_i^{R_m} |
        - \frac{1}{2}\nabla_1^2
        - \frac{1}{\lvert \bm{r}_1 - \bm{R}_1 \rvert}
        - \frac{1}{\lvert \bm{r}_1 - \bm{R}_2 \rvert}
        + \frac{1}{\lvert \bm{R}_1 -
            \bm{R}_2 \rvert} |  \psi_1}
    = E \Braket{\phi_i^{R_m} | \psi_1}.
\end{equation}

So the eigenvalue problem will be
\begin{equation}
    \det(T_{ij, mn} + N^E_{ij, mn} + N^N_{ij, mn} - S_{ij, mn} E ) = 0.
\end{equation}
where $T_{ij, mn} = \Braket{\phi_i^{R_m} | - \frac{1}{2}\nabla^2 | \phi_j^{R_n} }$,
$N^E_{ij, mn}=\Braket{\phi_i^{R_m} |
        -\frac{1}{\lvert \bm{r} - \bm{R}_1 \rvert } | \phi_j^{R_n} }
    + \Braket{\phi_i^{R_m} |
        -\frac{1}{\lvert \bm{r} - \bm{R}_2 \rvert } | \phi_j^{R_n} }$,
$N^N_{ij} = \Braket{\phi_i^{R_m} |
        \frac{1}{\lvert \bm{R}_1 - \bm{R}_2 \rvert } | \phi_j^{R_n} }$,
and
$S_{ij, mn} = \Braket{\phi_i^{R_m} | \phi_j^{R_n} }$
are all $8\times 8$ matrix,
where $i, j=1$,~$2$,~$3$,~$4$, and $m, n=1$,~$2$.

The total energy is
\begin{equation}
    \begin{aligned}
        E_{\text{total}} & = 2 E - \Braket{\psi | \frac{1}{\lvert \bm{R}_1 - \bm{R}_2 \rvert } | \psi}
                         & = 2 E - \frac{1}{\lvert \bm{R}_1 - \bm{R}_2 \rvert},
    \end{aligned}
\end{equation}
here the second term is calculated twice in eigenvalue, so subtract it.

When $\lvert \bm{R}_1 - \bm{R}_2 \rvert = \SI{1.020}{\bohr}$, the eigenvalue is \SI{-1.88723984843}{\hartree}.
See Fig. \ref{fig:pro3a} for reference.

% \begin{figure}[h]
%     \centering
%     \includegraphics[width=0.5\linewidth]{pro_3_a}
%     \caption{The total energy as a function of $\lvert \bm{R}_1 - \bm{R}_2 \rvert$.}
%     \label{fig:pro3a}
% \end{figure}

If I use the basis set from problem $1(a)$, the answers does not change.

\subsection*{(b)}

Now the Hamiltonian becomes
\begin{multline}
    H = - \frac{1}{2}\nabla_1^2
    - \frac{1}{2}\nabla_2^2
    - \frac{1}{\lvert \bm{r}_1 - \bm{R}_1 \rvert} - \frac{1}{\lvert \bm{r}_1 - \bm{R}_2 \rvert}\\
    - \frac{1}{\lvert \bm{r}_2 - \bm{R}_1 \rvert} - \frac{1}{\lvert \bm{r}_2 - \bm{R}_2 \rvert}
    + \frac{1}{\lvert \bm{R}_1 - \bm{R}_2 \rvert} + \frac{1}{\lvert \bm{r}_1 - \bm{r}_2 \rvert},
\end{multline}
and
\begin{multline}
    \Braket{\phi_i^{R_m} |
        - \frac{1}{2}\nabla_1^2
        - \frac{1}{\lvert \bm{r}_1 - \bm{R}_1 \rvert}
        - \frac{1}{\lvert \bm{r}_1 - \bm{R}_2 \rvert}
        + \frac{1}{\lvert \bm{R}_1 -
            \bm{R}_2 \rvert}
        + \int d\bm{r}_2 \frac{ \psi_2 ^\ast  (\bm{r}_2) \psi_2(\bm{r}_2)}{\lvert \bm{r}_1 - \bm{r}_2 \rvert}
        |  \psi_1}\\
    = E \Braket{\phi_i^{R_m} | \psi_1}.
\end{multline}
Similar as $2(b)$, the eigenvalue problem now becomes
\begin{equation}
    \det(T_{ij, mn} + N^E_{ij, mn} + N^N_{ij, mn} + H_{ij, mn} - S_{ij, mn} E ) = 0,
\end{equation}
where
\begin{multline}
    H_{ij, mn} = \sum_{\substack{p, q\\u, v}} 8\pi^2 c_{p, u}^\ast c_{q, v}
    \!\int_{0}^{\infty} r_1^2  \, dr_1
    \!\int_{0}^{\infty} r_2^2 \, dr_2
    \!\int_{0}^{\pi}
    \tfrac{\big( \phi_i^{R_m} \big)^\ast(r_1) \big(\phi_p^{R_u} \big)^\ast(r_2) %
        \phi_q^{R_v}(r_2)\phi_j^{R_n}(r_1)}{%
        \sqrt{r_1^2 + r_2^2 - 2 r_1 r_2 \cos\theta}} \sin\theta \, d\theta
\end{multline}
is the Hartree term.

The total energy is
\begin{equation}
    E_{\text{total}} = 2 E - \frac{1}{\lvert \bm{R}_1 - \bm{R}_2 \rvert} -
    \Braket{\psi | \frac{1}{\lvert \bm{r}_1 - \bm{r}_2 \rvert} | \psi },
\end{equation}
we need to subtract Hartree energy.

% \begin{figure}[h]
%     \subfigure[The total energy as a function of distance.]{
%         \label{fig:3-b:mini:subfig:a}   %% label for first subfigure
%         \begin{minipage}[b]{0.5\textwidth}
%             \centering
%             \includegraphics[width=\linewidth]{pro_3_b_1}
%         \end{minipage}}
%     \hfill
%     \subfigure[The density as a function of distance.]{
%         \label{fig:3-b:mini:subfig:b}   %% label for second subfigure
%         \begin{minipage}[b]{0.5\textwidth}
%             \centering
%             \includegraphics[width=\linewidth]{pro_3_b_2}
%         \end{minipage}}
%     \caption{Problem $3(b)$.}
%     \label{fig:3-b}   %% label for entire figure
% \end{figure}

See Fig. \ref{fig:3-b:mini:subfig:a} and Fig. \ref{fig:3-b:mini:subfig:b}
for reference. The minimum point is when $\lvert \bm{R}_1- \bm{R}_2 \rvert = \SI{1.3462}{\bohr}$(close
to experiment value \SI{1.401}{\bohr}), the
corresponding energy is \SI{-1.1129}{\hartree}.

\subsection*{(c)}

The Hamiltonian is still the same, but Schr\"{o}dinger equation will be
\begin{multline}\label{eq:psi1}
    \bigg(- \frac{1}{2}\nabla_1^2
    - \frac{1}{\lvert \bm{r}_1 - \bm{R}_1 \rvert}
    - \frac{1}{\lvert \bm{r}_1 - \bm{R}_2 \rvert}
    + \frac{1}{\lvert \bm{R}_1 -
        \bm{R}_2 \rvert}\\
    + \int d\bm{r}_2
    \frac{ \psi_2 ^\ast  (\bm{r}_2) \psi_2(\bm{r}_2) - \psi_2 ^\ast  (\bm{r}_2) \hat{P} \psi_2(\bm{r}_2)}%
    {\lvert \bm{r}_1 - \bm{r}_2 \rvert}
    \bigg) \psi_1(\bm{r}_1)
    = E_1  \psi_1(\bm{r}_1),
\end{multline}
and
\begin{multline}\label{eq:psi2}
    \bigg(- \frac{1}{2}\nabla_2^2
    - \frac{1}{\lvert \bm{r}_2 - \bm{R}_1 \rvert}
    - \frac{1}{\lvert \bm{r}_2 - \bm{R}_2 \rvert}
    + \frac{1}{\lvert \bm{R}_1 -
        \bm{R}_2 \rvert}\\
    + \int d\bm{r}_1
    \frac{ \psi_1 ^\ast  (\bm{r}_1) \psi_1(\bm{r}_1) - \psi_1 ^\ast  (\bm{r}_1) \hat{P} \psi_1(\bm{r}_1)}%
    {\lvert \bm{r}_1 - \bm{r}_2 \rvert}
    \bigg) \psi_2(\bm{r}_2)
    = E_2  \psi_2(\bm{r}_2),
\end{multline}
where $\hat{P}$ is the permutation operator.
So the eigenvalue problem will be
\begin{equation}
    \det(T_{ij, mn} + N^E_{ij, mn} + N^N_{ij, mn} + F_{ij, mn} - S_{ij, mn} E ) = 0.
\end{equation}
where
\begin{multline}
    F_{ij, mn} = \sum_{\substack{p, q\\u, v}} 8\pi^2 c_{p, u}^\ast c_{q, v}
    \!\int_{0}^{\infty} r_1^2  \, dr_1
    \!\int_{0}^{\infty} r_2^2 \, dr_2
    \!\int_{0}^{\pi}
    \frac{ \sin\theta \, d\theta}{ \mathsmaller{\sqrt{r_1^2 + r_2^2 - 2 r_1 r_2 \cos\theta}} }\\
    \cdot \Big(
    \big( \phi_i^{R_m} \big)^\ast(r_1) \big(\phi_p^{R_u} \big)^\ast(r_2) %
    \phi_q^{R_v}(r_2)\phi_j^{R_n}(r_1)\\
    -  \big( \phi_i^{R_m} \big)^\ast(r_1) \big(\phi_p^{R_u} \big)^\ast(r_2) %
    \phi_j^{R_n}(r_2)\phi_q^{R_v}(r_1)
    \Big)
\end{multline}
Now perform self-consistency loop, first assume a trial $\psi_2$, substitute into
$\eqref{eq:psi1}$, solve the coefficients of $\psi_1$, then substitute this $\psi_1$
into $\eqref{eq:psi2}$, solve the new coefficients of $\psi_2$, etc. Repeat this procedure
until the total energy converges.

\end{document}
